% =====================================================================
%  training_tables_molecular.tex
%  Complete training descriptions for the QM9 and GEOM-Drugs models.
%
%  EVERY value below was read from the model checkpoints themselves
%  (`model_args` / `args` / `student_cfg` / `hyper_parameters` /
%  `data_stats`), not from launcher scripts or notes, so the tables
%  describe the weights that actually produced the reported results.
%
%  Sources
%    QM9 base flow   MultiMFM/checkpoints/base_flow/model_step_100000.pt
%    QM9 MFM student MultiMFM/checkpoints/mfm_student/student_step_1000.pt
%    GEOM teacher    geom-drug-tabasco/validation_teacher/big_curve/
%                      ckpts/epoch=38-step=89310.ckpt
%    GEOM MFM student geom-drug-tabasco/outputs/mfm_student_geom_stage2_ddp/
%                      stage2_20260727-222154/student_step_10000.pt
%
%  These REPLACE tab:hyperparameters_multimfm (which describes the 2-D toy,
%  not QM9) and supply the GEOM table the paper currently lacks.
%
%  Generated 2026-09-25.
% =====================================================================

\documentclass[11pt]{article}
\usepackage[margin=1in]{geometry}
\usepackage{booktabs,amsmath,amssymb,xcolor}
\newcommand{\naa}{\textcolor{gray}{--}}
\title{Molecular training details: QM9 and GEOM-Drugs}
\date{2026-09-25}
\begin{document}
\maketitle

\section{QM9}

\begin{table}[h]\centering
\caption{QM9 training. The base flow model is trained unconditionally; the
multiMFM student is then distilled from it. Both are read from the released
checkpoints. Note the student is a \emph{fine-tune}: it is initialised from an
earlier student checkpoint rather than from scratch or from the teacher, and the
reported weights are taken at update 1{,}000 of a 4{,}000-update cosine schedule.}
\label{tab:qm9-training}
\begin{tabular}{lrr}
\toprule
\textbf{Hyperparameter} & \textbf{Base flow} & \textbf{multiMFM student} \\
\midrule
\multicolumn{3}{l}{\emph{Data}} \\
Dataset                        & QM9 (\texttt{yairschiff/qm9}) & QM9 \\
Partition                      & \texttt{train\_flow} & \texttt{train\_flow} \\
Molecules                      & $50{,}000$ & $40{,}000$ \\
Hydrogens                      & excluded & excluded \\
Maximum heavy atoms            & $9$ & $9$ \\
Atom classes $K$               & $9$ & $9$ \\
Coordinate normalizer          & $1.0$ & $1.0$ \\
\midrule
\multicolumn{3}{l}{\emph{Architecture}} \\
Hidden dimension               & $64$ & $64$ \\
Transformer layers             & $4$  & $4$ \\
Attention heads                & $4$  & $4$ \\
Encoder depth                  & \naa & $4$ \\
Block conditioning             & \naa & additive \\
Velocity parameterization      & \naa & endpoint \\
Total parameters               & $338{,}121$ & $363{,}465$ \\
\midrule
\multicolumn{3}{l}{\emph{Probability path}} \\
Atom process                   & \multicolumn{2}{c}{Gaussian prior $\to$ simplex (\texttt{dfm\_gaussian})} \\
$\beta$ schedule               & \multicolumn{2}{c}{linear, $\beta(t)=t$} \\
Coordinate path                & \multicolumn{2}{c}{linear interpolant, zero-COM} \\
Sampling schedule              & linear & linear \\
\midrule
\multicolumn{3}{l}{\emph{Optimization}} \\
Optimizer                      & Adam & Adam \\
Learning rate                  & $2\times10^{-3}$ & $1\times10^{-6}$ \\
Warmup steps                   & \naa & $200$ \\
Schedule                       & constant & cosine (horizon $4{,}000$) \\
Batch size                     & $128$ & $4$ \\
Training steps                 & $100{,}000$ & $1{,}000$ (of $4{,}000$) \\
EMA decay                      & \naa & $0.999$ \\
Gradient clipping ($\ell_2$)   & \naa & $1.0$ \\
Seed                           & $0$ & $44$ \\
\midrule
\multicolumn{3}{l}{\emph{Distillation objective (student only)}} \\
Stage                          & \naa & 1 (diagonal) \\
Diagonal weight                & \naa & $1.0$ \\
Coordinate / atom weights      & \naa & $8.0\,/\,0.5$ \\
Diagonal direction weight      & \naa & $5.0$ \\
Sensitivity weight (L2)        & \naa & $0.1$ \\
Sensitivity weight (direction) & \naa & $0.25$ \\
Sensitivity $\varepsilon$      & \naa & $0.03$ \\
Consistency (ESD) weight       & \naa & \textbf{not used at stage 1} \\
Maximum training time $t_{\max}$ & \naa & $0.95$ \\
$t_{\mathrm{cond}}$ grid       & \naa & $\{0.1,0.3,0.5,0.7,0.9\}$, jitter $0.04$ \\
Conditioning states            & \naa & teacher-generated ($128$ steps) \\
Teacher posterior steps        & \naa & $16$ \\
Initialised from               & random & earlier student checkpoint \\
\bottomrule
\end{tabular}
\end{table}

\paragraph{Two disclosures the current text omits.}
\begin{enumerate}
\item \textbf{The QM9 student is trained with the diagonal and sensitivity terms
only; it carries no ESD term.} The main-text objective
$\mathcal{L}=\mathcal{L}_{\mathrm{diag}}+\mathcal{L}_{\mathrm{con}}$ and
Algorithm~1 therefore do not describe this model. State which models use
consistency training and which do not.
\item \textbf{The student is a fine-tune of an earlier checkpoint}
($\texttt{lr}=10^{-6}$, batch $4$, $1{,}000$ updates), so its cumulative training
budget exceeds what the table shows. The parent checkpoint was lost to a disk
purge, so the full lineage cannot be reconstructed; say so rather than implying
the model was trained in $1{,}000$ updates.
\end{enumerate}

\paragraph{Definition of the sensitivity term.} Let $x_{\mathrm{cond}}$ be the
conditioning state and $\delta$ a masked, zero-COM perturbation rescaled so that
its per-sample RMS equals $\varepsilon$. With $v_\theta$ the student velocity and
$v_{\mathrm{teacher}}$ the teacher's GLASS velocity,
\[
\mathcal{L}_{\mathrm{sens}}
=\Big\|\,
\tfrac{1}{\varepsilon}\big(v_\theta(x_{\mathrm{cond}}{+}\delta)-v_\theta(x_{\mathrm{cond}})\big)
-\tfrac{1}{\varepsilon}\big(v_{\mathrm{teacher}}(x_{\mathrm{cond}}{+}\delta)-v_{\mathrm{teacher}}(x_{\mathrm{cond}})\big)
\Big\|^2 ,
\]
plus a directional term $1-\cos(\cdot,\cdot)$ on the same two differences. It is a
central finite difference with respect to the conditioning \emph{state} --- not the
conditioning time, and not a forward-mode JVP despite the internal name.

\section{GEOM-Drugs}

\begin{table}[h]\centering
\caption{GEOM-Drugs training. The teacher is trained unconditionally on
GEOM-Drugs; the multiMFM student is distilled from its EMA weights. Values are
read from the checkpoints. The student is the capped-gap stage-2 model used for
the reported fingerprint-guidance results.}
\label{tab:geom-training}
\begin{tabular}{lrr}
\toprule
\textbf{Hyperparameter} & \textbf{Teacher} & \textbf{multiMFM student} \\
\midrule
\multicolumn{3}{l}{\emph{Data}} \\
Dataset                        & GEOM-Drugs & GEOM-Drugs \\
Molecules                      & $1{,}172{,}182$ & $150{,}000$ \\
Hydrogens                      & excluded & excluded \\
Maximum heavy atoms            & $71$ & $71$ \\
Atom classes $K$               & $9$ & $9$ \\
Atom vocabulary                & \multicolumn{2}{c}{C, N, O, F, S, Cl, Br, I, $*$} \\
Coordinate normalizer          & $2.0$ & $2.0$ \\
\midrule
\multicolumn{3}{l}{\emph{Architecture}} \\
Hidden dimension               & $256$ & $256$ \\
Transformer layers             & $16$  & $16$ \\
Attention heads                & $8$   & $8$ \\
Encoder depth                  & \naa  & $16$ \\
Block conditioning             & \naa  & additive \\
Velocity parameterization      & \naa  & endpoint \\
Total parameters               & $14{,}795{,}529$ & $15{,}191{,}817$ \\
\midrule
\multicolumn{3}{l}{\emph{Probability path}} \\
Atom process                   & \multicolumn{2}{c}{Gaussian prior $\to$ simplex (\texttt{dfm\_gaussian})} \\
$\beta$ schedule               & \multicolumn{2}{c}{linear, $\beta(t)=t$} \\
Sampling schedule              & log & log \\
Random augmentations           & $7$ ($8\times$ effective batch) & inherited \\
\midrule
\multicolumn{3}{l}{\emph{Optimization}} \\
Optimizer                      & Adam ($\beta=[0.9,0.995]$) & Adam \\
Learning rate                  & $2\times10^{-3}$ & $1\times10^{-4}$ \\
Weight decay                   & $0.0$ & $0.0$ \\
Warmup steps                   & \naa & $500$ \\
Batch size (per GPU)           & $256$ & $96$ \\
GPUs                           & $2$ (DDP) & $4$ (DDP) \\
Global effective batch         & $4{,}096$ ($8\times$ aug.) & $384$ \\
Training steps                 & $89{,}310$ (epoch $38$) & $10{,}000$ (of $30{,}000$) \\
EMA decay                      & $0.999$ & $0.995$ \\
Seed                           & \naa & $0$ \\
\midrule
\multicolumn{3}{l}{\emph{Distillation objective (student only)}} \\
Stage                          & \naa & 2 (diagonal $+$ ESD) \\
Diagonal weight                & \naa & $2.0$ \\
Consistency (ESD) weight       & \naa & $1.0$ \\
Coordinate / atom weights      & \naa & $1.0\,/\,1.0$ \\
Sensitivity weights            & \naa & $0$ (disabled) \\
ESD target clip                & \naa & $10.0$ \\
ESD delta clip                 & \naa & $8.0$ \\
Maximum jump $|u-s|$           & \naa & $0.15$ \\
Anneal fraction                & \naa & $0.2$ \\
Maximum training time $t_{\max}$ & \naa & $0.95$ \\
$t_{\mathrm{cond}}$ power / zero rate & \naa & $2.0$ / $0.1$ \\
Teacher posterior steps        & \naa & $64$ \\
Teacher weights used           & \naa & EMA \\
Initialised from               & random & earlier stage-2 checkpoint \\
\bottomrule
\end{tabular}
\end{table}

\paragraph{Disclosures.}
\begin{enumerate}
\item \textbf{The GEOM student is also a resume}, from an earlier stage-2
checkpoint, so its $10{,}000$ updates are not its full budget.
\item \textbf{The teacher is selected at epoch 38 by EMA validation loss}, not at
the end of training; several longer runs exist. Name the selection criterion.
\item \textbf{The maximum-jump cap of $0.15$ is load-bearing}: without it the ESD
term diverges once the annealed gap reaches ${\sim}0.20$. This is worth one
sentence, since it explains the clipping hyperparameters above.
\end{enumerate}

\section{Sampling settings used for the reported results}

\begin{table}[h]\centering
\caption{Inference-time settings. These are the configurations behind the
reported numbers, not defaults --- the released scripts' argparse defaults differ
and will not reproduce the results on their own.}
\label{tab:molecular-sampling}
\begin{tabular}{lrr}
\toprule
 & \textbf{QM9} & \textbf{GEOM-Drugs} \\
\midrule
Outer integration steps        & $128$ & $100$ \\
Inner flow-map steps           & $4$ (diagonal) & $10$ \\
Posterior (Monte Carlo) samples & $32$ & $4$ \\
Guidance applied every         & step & $2$ steps \\
Guidance window $t$            & $[0.05,\,0.95]$ & $[0.1,\,0.9]$ \\
Soft atom temperature          & $0.25$ & $0.25$ \\
Guidance strength $\mu$        & per property & $0.5$ \\
Reward scale $\lambda$         & per property & $0.02$ \\
Teacher weights                & final & EMA \\
\bottomrule
\end{tabular}
\end{table}

\paragraph{Reward.} For a target property value $y^{\star}$ and differentiable
guide predictor $\hat y$, the terminal reward is
$r(x_1)=-\big(\hat y(x_1)-y^{\star}\big)^2$, and the value function is
$V_t(x)=\log\mathbb{E}\!\left[e^{\lambda r(X_1)}\mid X_t=x\right]$, estimated with
the multiMFM posterior samples above. Atom types enter the guide through a
softmax relaxation at temperature $0.25$; all \emph{reported} metrics use the hard
\texttt{argmax} readout, so the differentiated objective and the evaluated
objective are not the same function --- worth stating once.

\end{document}
